%%%%%%%%%%  scan idx 42  |  folio 29  %%%%%%%%%%
le
polynôme $Q_i$ étant isobare de poids $i$. De plus, $x*y$ doit être distributif par rapport à l'«~addition~» $xy$ et pour des facteurs «~linéaires~» doit se réduire à ce qui suit~:
\begin{equation}\tag{1.17}
[1,1+x^1][1,1+y^1]=[1,1+x^1+y^1].
\end{equation}
Les polynômes $Q_i$ s'évaluent par des calculs de fonctions symétriques élémentaires (il s'agit du calcul des classes de Chern d'un produit tensoriel de fibrés vectoriels à l'aide des classes de Chern $x^i$ et $y^i$ des facteurs et de leurs dimensions respectives $m$ et $n$). De (1.17) on déduit~:
\begin{equation}\tag{1.17 bis}
(1+x^1)*(1+y^1)=(1+x^1+y^1)/(1+x^1)(1+y^1),
\end{equation}
et un produit $x*y$ quelconque se calcule terme à terme en décomposant «~formellement~» $x$ en un produit $\prod_{i=1}^n(1+\alpha_i)$ et $y$ en un produit $\prod_{i=1}^n(1+\beta_i)$, où les $\alpha_i$ et les $\beta_i$ sont de degré $1$.

De cette manière, $\widetilde A$ devient un anneau commutatif augmenté, avec une unité $[1,1]$, et on peut l'interpréter comme l'anneau obtenu par adjonction d'une unité à l'anneau $1+\Ahp$. De plus, sur $\widetilde A$, on définit une \emph{filtration} en posant $\widetilde A^0=\widetilde A$ et en notant $\widetilde A^n$ l'ensemble des $[0,x]$ avec $x^i=0$ pour $1\leq i<n$~; cette filtration est compatible avec le produit, car on a
\begin{equation}\tag{1.18}
\left(1+\sum_{i\geq m}x^i\right)*\left(1+\sum_{j\geq n}y^j\right)
=1-\frac{(m+n-1)!}{(m-1)!(n-1)!}x^m y^n+\cdots,
\end{equation}
les termes non écrits étant de degré $>m+n$. Cette formule montre que l'anneau gradué $G(\widetilde A)$ associé à $\widetilde A$ s'identifie à $A$ au point de vue additif, mais non au point de vue multiplicatif~; on peut dire seulement que l'application $\eta'$ de $A$ dans $G(\widetilde A)$ définie par
